\documentclass[border = 1mm]{standalone}
% Caption:
% — Torus.
% Parametric surface defined by the points in space (x,y,z)
% such that
% 	x = (R + r cos(v)) * cos(u)
%	y = (R + r cos(v)) * sin(u)
%	z = r sin(v)
% for R and r positive numbers, and u,v in [0, 2*pi).
% Equivalently, the torus is the implicit surface
%
%   (sqrt(x^2+y^2)-R)^2 + z^2 = r^2
%
% obtained by revolving a circle of radius r centered
% at distance R from the z-axis.
%
% Autor:
% R. Sánchez-Martínez, Jun 2026.
%
% Figure published at:
% instagram.com/ars_mathematica
% Code available at:
% github.com/ro-smtz/ars-mathematica-content/tree/main/figures
%
% Generated with TikZ and compiled preferable with XeLaTeX.

%
\usepackage[utf8]{inputenc}
\usepackage[T1]{fontenc}
\usepackage{microtype}
%
\usepackage{amsmath,amssymb,amsfonts}%
\usepackage{xcolor}
\usepackage{tikz}
\usetikzlibrary{
	babel,
	math,calc
}
\usepackage{pgfplots}
\pgfplotsset{compat = 1.18}

\usepackage{fontspec}
\usepackage[default]{fontsetup}

\begin{document}
% -- Figure
\begin{tikzpicture}[
	line join = round,
	line cap  = round,
	line width = 0.62pt,
]

	% Main box 
	\node[
		inner sep = 0pt,
		minimum width  = 6.2cm,
		minimum height = 6.2cm,
		name = N,
	] at (0,0) {};
	
	% Clipping mask, it deletes everything outside the box (N),
	% useful for obtaining well-defined-size drawings
	\clip (N.north west) rectangle (N.south east);

	% Geometric parameters:
	%   R : major radius
	%       Distance from the center of the torus
	%       to the center of the tube.
	%
	%   r : minor radius
	%       Radius of the tube itself.
	\tikzmath{
		\R = 1.4;
		\r = 0.8;
	};

	% ---
	% Main drawing
	\begin{axis}[
		% position
		at = {(N)},
		anchor = center,
		shift = {(2mm,1mm)},
		% size
		width = 7cm,
		view = {56}{20},
		unit vector ratio = 1 1 1,
		%
		xmin = -2, xmax = 2,
  		ymin = -2, ymax = 2, 
  		zmin = -1.5, zmax = 2,
  		%
  		enlarge x limits,
  		enlarge y limits,
  		%
  		grid,
		% 
		xlabel = {$x$},
 		ylabel = {$y$},
 		zlabel = {$z$},
 		%
 		xticklabels = \empty,
  		yticklabels = \empty,
  		zticklabels = \empty,
 		%
 		trig format plots = rad,
 		% Two-point grayscale colormap used to shade the surface
 		%
		% The first entry assigns white (1) to the beginning
		% of the colormap, while the second assigns dark gray
		% (0.22) to its end — 0 corresponds to black.
		% PGFPlots interpolates smoothly between both values.
 		colormap = {blackwhite}{
			gray(0cm) = (1);
			gray(1cm) = (0.22)
		},
	]
	
	% The torus is plotted in two separate patches.
	%
	% Splitting the surface into upper and lower halves
	% improves visibility of the interior and allows the
	% generatrix circle to remain visible through the
	% transparent surface.
	
		% lower half surface
		\addplot3[
			surf,
			z buffer = sort,
			domain = 0*pi:1.5*pi, samples = 20,
			domain y = 1*pi:2*pi, samples y = 10,
			opacity = 0.75,
		]
		(
			{(\R + \r*cos(y))*cos(x)},
			{(\R + \r*cos(y))*sin(x)},
			{\r*sin(y)}
		);
		\addlegendentry{$S$}
  		
  		% inner circle, generatrix of the torus
  		\addplot3[
			semithick, black,
  			domain = 0:2*pi, 
  			samples = 120, samples y = 0,
  		]
  		(
  			{\R*cos(x)},
  			{\R*sin(x)},
  			{0}	
  		);
  		
  		% ppper half surface
		\addplot3[
			surf,
			z buffer = sort,
			domain = 0*pi:1.5*pi, samples = 20,
			domain y = 0*pi:1*pi, samples y = 10,
			opacity = 0.75,
		]
		(
			{(\R + \r*cos(y))*cos(x)},
			{(\R + \r*cos(y))*sin(x)},
			{\r*sin(y)}
		);
  		
  		% generating circle (meridian section)
  		\addplot3[
			semithick, black,
  			y domain = 0:2*pi, 
  			samples = 0, samples y = 120,
  		]
  		(
  			{(\R + \r*cos(y))*cos(1.5*pi)},
  			{(\R + \r*cos(y))*sin(1.5*pi)},
  			{\r*sin(y)}	
  		);

		% Auxiliary points used to construct dimension-like
		% annotations for the major radius R and minor radius r.
		\coordinate (O) at (axis cs: 0, 0    ,0);
		\coordinate (P1) at (axis cs: 0,-\R   ,0);
		\coordinate (P2) at (axis cs: 0,-\R-\r,0);
		
		% (axis direction cs: x,y,z) shifts points along the local
		% coordinate directions of the 3D axis rather than
		% in the surrounding TikZ canvas coordinates.
		\path (O)++(axis direction cs: 0,0,1.7) coordinate (O-up);
		\path (P1)++(axis direction cs: 0,0,1.7) coordinate (P1-up);
		\path (P2)++(axis direction cs: 0,0,1.7) coordinate (P2-up);
			
	\end{axis}
	% ---
	
	% radii annotations
	\draw[gray!75, shorten <= 1mm]
		(O) to
			coordinate[pos = 0.9] (O-aux)
		(O-up);
	\draw[gray!75, shorten <= 1mm]
		(P1) to
			coordinate[pos = 0.9] (P1-aux)
		(P1-up);
	\draw[gray!75, shorten <= 2mm]
		(P2) to
			coordinate[pos = 0.9] (P2-aux)
		(P2-up);
		
	\draw[gray!75]
		(P2-aux) to 
			node[shift = {(0,0.25)}, black] {$r$}
		(P1-aux) to
			node[shift = {(-0.05,0.25)}, black] {$R$}
		(O-aux);
		
	\fill (O) circle (1.1pt);
	\fill (P1) circle (1.1pt);
	
\end{tikzpicture}
\end{document}